\chapter{Conclusion}

GW150914 opened to the door to GW astronomy. Almost everything humanity currently knows about the universe was learned through electromagnetic observations. GWs are generated from a different force and offer an entirely new approach to studying the universe. It is not an exaggeration to compare the discovery of GW150914 to Galileo's first examination of the sky with a telescope. One or two hundred years from now, humanity's understanding of the universe may be just as fundamentally dependent on GW astronomy as it currently is on electromagnetic astronomy.

Compact binary systems will likely dominate GW observations in the near future, but new frequency bands and GW sources will become increasingly important as detectors and sensitivities improve. CCSNe represent one of the most promising and mysterious sources of GWs yet to be observed. These GWs originate in the core of a collapsing star and carry important information about processes and mechanisms that are not fully understood. An observation of a GW from a CCSN could provide clues to the nuclear equation of state and could elucidate the exact physics behind the spread of heavy metals throughout the universe. These processes, while seemingly remote and obscure, are essential to all known life.

Spectrograms offer a robust way to approach and study GW emission from a CCSN signal. The spectrogram version of SMEE presented in this dissertation does not use unreliable phase data. Instead, its statistics depend entirely on the power, frequency, and time of GW emission. This is reflected in the reduction of the number of PCs needed from those in previous studies that used the time series waveforms. The time-frequency path, inherent to a spectrogram, also allows the study and analysis of specific waveform features. This results in a robust and sensitive tool to perform the difficult task of parameter estimation of a GW signal from a CCSN. The small number of simulated waveforms is still a concern, as is the fact that most simulations are ended prematurely. As more simulations continue to be released they will be incorporated into SMEE's analysis.

The work in this dissertation contains the first study of CCSN waveforms in future detector networks. SMEE's performance in future detectors was simulated and mapped out as realistically as possible with recolored Advanced LIGO data. The results suggest that SMEE should be able to classify CCSN waveforms from well outside of the Milky Way in future detector configurations. Third generation detectors should be able to resolve the explosion mechanism and waveform features for most CCSN out to beyond the Small and Large Magellanic Clouds (50 - 60\,kpc), as well as for the dozen or so satellite galaxies in between~\cite{karachentsev2004catalog, belokurov2007cats}. A fraction of magnetorotational waveforms should be classifiable all the way out to the Andromeda galaxy (790\,kpc) in third generation detectors, while neutrino model performance falls off significantly beyond 150\,kpc. There are a total of 28 known Galaxies, mostly satellite or dwarf, within the range of 150\,kpc, and 55 within the range of 790\,kpc~\cite{karachentsev2004catalog, belokurov2007cats}. While CCSN are rare (1 - 2 per galaxy per century~\cite{1991ARA&A..29..363V, 1993A&A...273..383C, Alexeyev2002}), these results suggest that in third generation configurations a detection and accurate classification are both very plausible.

This dissertation represents the first to attempt to identify specific features associated with g-modes and SASI in detected GW signals from CCSNe. If more common features are found in future CCSN simulations, they could be incorporated into SMEE's analysis. The observation of a gravitational wave from a CCSN will be an important moment in astrophysics, and with a tool such as SMEE it will be possible to learn about the source and the relevant underlying physics of the CCSN explosion.

